\honest{The Allais Paradox}{Eliezer Yudkowsky}{2008}

Choose: \$24,000 for certain, or a 33/34 chance of \$27,000. Now choose again: a 34\% chance of \$24,000, or a 33\% chance of \$27,000. Most people take the sure thing the first time and the bigger prize the second, and in within-subject comparisons a majority do both. \nb{Kahneman and Tversky (1979), whom the post lists, found 61 per cent making both choices in their version. The post's numbers are a modified form that breaks the same axiom.}

This is a problem because each second gamble is the first played with probability 0.34. The independence axiom says that mixing two gambles with the same third gamble should not reverse your preference between them. So no utility function, whatever its shape, lets you make both choices and maximize expected utility, and the diminishing value of money has nothing to do with it. \nb{The algebra is right. Theories that drop independence but keep a consistent ranking of gambles, such as rank-dependent utility, can represent both choices.}

Allais ``initially defended'' his subjects' preferences and took the experiment to expose a flaw in the theory, not in human psychology. ``This was 1953, after all.'' \nb{Allais still held the view in 1979, when he published a 245-page rejoinder extending it.} Then, in sarcasm: ``How naive, how foolish, how simplistic is Bayesian decision theory.'' Surely certainty should count for something; you can feel it. I call this ``naive philosophical realism'': intuitions tell you about your own mind, and only indirectly about which choice is wise. \nb{The conclusion, that the Allais choices are a mistake, was Savage's in 1954. Savage reported that he had first made the Allais choices himself and then changed his mind. On the other side, Buchak (2013) argues that a rational agent may weigh risk itself, so that the Allais choices are ``rationally permissible,'' in the Stanford Encyclopedia's summary.}

Why fret about elegance? Because ``There is always a price for leaving the Bayesian Way.'' The price is a money pump. At noon I roll a hundred-sided die; above 34, the game ends. Otherwise a switch decides whether I pay you \$24,000, or \$27,000 unless a 34-sided die shows 34. Before the roll you pay me a penny to set the switch to the gamble. The die shows 12, and you pay me a penny to set it back. ``I have taken your two cents on the subject.'' \nb{That works on an agent who does not look ahead. An agent who plans ahead would not pay the first penny, and would end with the choice it likes less. Defenders of Allais's view reply that an agent may make a plan and keep to it. And when Kahneman and Tversky described a game of this shape in two stages, most subjects chose the sure option from the start.}

I think the same failure to discount small probabilities explains the lottery. \nb{Kahneman and Tversky say the same: overweighting of low probabilities ``may contribute to the attractiveness of both insurance and gambling.''}
